\chapter{Solutions}

\section{Exercices essentiels}\label{II-exercices-essentiels}

\runinhead{Approximation par une
constante}\phantomsection\addcontentsline{toc}{subsubsection}{Approximation
par une constante}\label{II-answer-exo-approx}
On a
\[\mathbb{E}((X-a)^2) = \mathbb{E}((X-\mathbb{E}(X) + \mathbb{E}(X) - a)^2 = \mathbb{E}((X - \mathbb{E}(X))^2 ) + (\mathbb{E}(X) - a)^2\]
le terme de gauche ne dépend pas de \(a\) tandis que celui de droite
s'annule en \(a = \mathbb{E}(X)\).

\runinhead{Calcul d'une
covariance}\phantomsection\addcontentsline{toc}{subsubsection}{Calcul
d'une covariance}\label{II-answer-cov-unif}
On a \begin{align*}
\text{Cov}(U,V) &= \mathbb{E}((U-\mathbb{E}(U))(V-\mathbb{E}(V)))\\
          &= \mathbb{E}(UV - \frac{U}{2} - \frac{U}{2} - \mathbb{E}(U)\mathbb{E}(V))\\
          &= \mathbb{E}(U(1-U)) - 1/4 - 1/4 +1/4 \text{ par linéarité et d'après la solution ci-dessous}\\
          &= 1/2 - 1/3 -1/4 = -1/12
\end{align*} et
\[\rho(U,V) = \frac{\text{Cov}(U,V)}{\sqrt{\mathbb{V}(U)}\sqrt{\mathbb{V}(V)}} = -1\]

\runinhead{Coefficient de
corrélation}\phantomsection\addcontentsline{toc}{subsubsection}{Coefficient
de corrélation}\label{II-answer-CS-corr}
On déduit de l'inégalité de
\protect\hyperref[II-CS]{Cauchy-Schwarz (p.~\pageref*{II-CS})} que
\[|\rho(X,Y)| = \left|\frac{\mathbb{E}((X-\mathbb{E}(X))(X-\mathbb{E}(Y)))}{\sqrt{\mathbb{V}(X)\mathbb{V}(Y)}} \right| \leq \frac{\sqrt{\mathbb{E}((X-\mathbb{E}(X))^2)\mathbb{E}((X-\mathbb{E}(Y))^2)}}{\sqrt{\mathbb{V}(X)\mathbb{V}(Y)}} =1\]

\runinhead{Moments d'une variable aléatoire de loi uniforme sur
\([a,b]\)}\phantomsection\addcontentsline{toc}{subsubsection}{Moments
d'une variable aléatoire de loi uniforme sur
\([a,b]\)}\label{II-answer-moments-unif}
Soit \(X \sim \mathcal{U}_{[a,b]}\). Son espérance vaut
\[\mathbb{E}(X) = \int_a^b \frac{x}{b-a} dx = \frac{a+b}{2}\] et puisque
\[\mathbb{E}(X^2) = \int_a^b \frac{x^2}{b-a} dx = \frac{a^2 + ab +b^2}{3},\]
alors sa variance vaut
\[\mathbb{V}(X) = \mathbb{E}(X^2) - \mathbb{E}(X)^2 = \frac{(b-a)^2}{12}.\]

\runinhead{Moments d'une v.a. de loi
exponentielle}\phantomsection\addcontentsline{toc}{subsubsection}{Moments
d'une v.a. de loi exponentielle}\label{II-answer-moments-expo}
Soit \(X \sim \mathcal{E}(\theta)\), alors
\[\mathbb{E}(X) = \int_0^{+\infty} x \theta \exp(-\theta x) dx = [\frac{x}{\theta}\exp(-\theta x)]_0^{+\infty} + \int_0^{+\infty} \exp(-\theta x) dx = \frac{1}{\theta}\]

\runinhead{Moments d'une v.a. de loi
gamma}\phantomsection\addcontentsline{toc}{subsubsection}{Moments d'une
v.a. de loi gamma}\label{II-answer-moments-gamma}
Soit \(X\sim \Gamma(\alpha,\beta)\). On a \begin{align*}
\mathbb{E}(X) &= \frac{1}{\Gamma(\alpha)}\int_0^{+\infty} \theta^\alpha x^\alpha e^{-\theta x} dx\\
        &= \frac{1}{\Gamma(\alpha)}\int_0^{+\infty} \frac{1}{\theta} u^\alpha e^{-u} du \text{ par le changement de variable } u = \theta x\\
        &= \frac{\Gamma(\alpha+1)}{\Gamma(\alpha)} \frac{1}{\theta} = \frac{\alpha}{\theta}
\end{align*} puis \begin{align*}
\mathbb{E}(X^2) &= \frac{1}{\Gamma(\alpha)}\int_0^{+\infty} \theta^\alpha x^{\alpha+1} e^{-\theta x} dx\\
        &= \frac{1}{\Gamma(\alpha)}\int_0^{+\infty} \frac{1}{\theta^2} u^{\alpha+1} e^{-u} du \text{ par le changement de variable } u = \theta x\\
        &= \frac{\Gamma(\alpha+2)}{\Gamma(\alpha)} \frac{1}{\theta^2} = \frac{\alpha(\alpha+1)}{\theta^2}
\end{align*} d'où
\[\mathbb{V}(X) = \mathbb{E}(X^2) - \mathbb{E}(X)^2 = \frac{\alpha}{\theta^2}\]

\runinhead{Moments d'une v.a.
gaussienne}\phantomsection\addcontentsline{toc}{subsubsection}{Moments
d'une v.a. gaussienne}\label{II-answer-moments-gauss}
Pour s'assurer qu'il s'agit bien d'une densité, on remarque que
\(I = \int_{-\infty}^{+\infty} f(x) dx\) vérifie
\[ I^2 = \int_{-\infty}^{+\infty} \int_{-\infty}^{+\infty} f(x) f(y) dxdy = \int_{0}^{2\pi} d\theta\int_{0}^{+\infty}\frac{1}{2\pi} r e^{-r^2/2}dr \]
(en passant en coordonnées polaires dans l'intégrale double). Le calcul
est ensuite aisé et on obtient \(I = 1\).

Pour les moments, on fait d'abord le calcul dans le cas centré réduit
(\(\mu = 0\) et \(\sigma^2 =1\)) puis on s'y ramène par le changement de
variable \(x \mapsto \frac{x-\mu}{\sigma}\). Soit donc
\(X \sim \mathcal{N}(0,1)\), on a \begin{align*}
\mathbb{E}(X) &= \int_{-\infty}^{+\infty} \frac{x}{\sqrt{2\pi}} e^{-x^2/2} dx\\
         &= \left[-\frac{1}{\sqrt{2\pi}} e^{-x^2/2} \right]_{-\infty}^{+\infty} = 0.
\end{align*} et \begin{align*}
\mathbb{V}(X) = \mathbb{E}(X^2) &= \int_{-\infty}^{+\infty} \frac{x^2}{\sqrt{2\pi}} e^{-x^2/2} dx\\
         &= \left[-\frac{x}{\sqrt{2\pi}} e^{-x^2/2} \right]_{-\infty}^{+\infty} +\int_{-\infty}^{+\infty} \frac{1}{\sqrt{2\pi}} e^{-x^2/2} dx \\
         &= 0 + 1.
\end{align*}

\runinhead{\emph{Loi de
Cauchy}}\phantomsection\addcontentsline{toc}{subsubsection}{\emph{Loi de
Cauchy}}\label{II-answer-exo-cauchy}
L'angle \(\theta\) est une variable aléatoire uniforme sur
\([-\pi/2,\pi/2]\), de densité
\(g(\theta) = \frac{1}{\pi}1_{[ -\pi/2,\pi/2 ]}(\theta)\). L'abscisse
\(X\) est donnée par \(X = \tan \theta\), c'est donc une variable
aléatoire, de fonction de répartition \begin{align*}        
F(x) & = \mathbb{P}(X \leq x) \\
     & = \mathbb{P}(\theta \leq \arctan x) \\
     & = \int_{-\infty}^{\arctan x} \frac{1}{\pi}1_{[ -\pi/2,\pi/2 ]}(\theta) d\theta \\
     & = \frac{1}{\pi} \arctan x + \frac{1}{2}.
\end{align*}

\(F\) est de classe \(C^1\) de dérivée
\[f(x) = \frac{1}{\pi(1+x^2)}, x \in \mathbb{R}\]

C'est la densité de la loi de Cauchy. Une variable aléatoire \(X\) de
loi de Cauchy n'admet pas d'espérance. En effet,
\(\frac{x}{\pi(1+x^2)} \sim_{x\to \infty} \frac{1}{x}\) n'est pas
intégrable.

\runinhead{Calcul de l'aire d'un
triangle}\phantomsection\addcontentsline{toc}{subsubsection}{Calcul de
l'aire d'un triangle}\label{II-answer-triangle}
Si la fonction \(1_D\) est intégrable, alors
\protect\hyperref[II-fubiniproba]{le théorème de Fubini est
applicable (p.~\pageref*{II-fubiniproba})}. On a donc \[
a(T) = \int_{\mathbb{R}}\left[ \int_{\mathbb{R}} 1_T(x, y) \, dx \right] \, dy.
\] Or, si \(0 \leq y \leq 1\), \[
1_T(x, y) = \left| 
\begin{array}{rl}
1 & \mbox{si $0 \leq x \leq 1-y$,} \\
0 & \mbox{sinon.}
\end{array}
\right.
\] et dans le cas contraire, \(1_T(x, y) = 0\). On a donc \[
a(T) = \int_0^1\left[ \int_0^{1-y} dx \right] \, dy
=\int_0^1 (1 - y) \, dy= \left[y - \frac{y^2}{2} \right]_0^1 = \frac{1}{2}.
\]

\runinhead{Contre-exemple}\phantomsection\addcontentsline{toc}{subsubsection}{Contre-exemple}\label{II-answer-fubini-counter-example}
Pour tout \(y \in \left]0,1\right]\) (et donc presque tout
\(y \in [0, 1]\)), on a \[
\int_0^1 \frac{x^2 - y^2}{(x^2+y^2)^2} \, dx 
= 
\left[x \mapsto - \frac{x}{x^2+y^2} \right]_0^1
= 
- \frac{1}{1+y^2},
\] par conséquent \[
\int_0^1 \left[ \int_0^1 \frac{x^2 - y^2}{(x^2+y^2)^2} \, dx \right] \, dy
= - \int_0^1 \frac{dy}{1+y^2} \\
= - [y \mapsto \arctan y]_0^1
= - \frac{\pi}{4}.
\] En exploitant la relation \[
\frac{x^2 - y^2}{(x^2+y^2)^2} = \frac{\partial}{\partial y} \left(\frac{y}{x^2+y^2}\right) 
\] on établit de façon similaire que \[
\int_0^1 \left[ \int_0^1 \frac{x^2 - y^2}{(x^2+y^2)^2} \, dy \right] \, dx
= \frac{\pi}{4}.
\]

Or, si \protect\hyperref[II-fubiniproba]{le théorème de
Fubini (p.~\pageref*{II-fubiniproba})} était applicable, on pourrait
intervertir l'ordre d'intégration des variables sans changer le
résultat. Comme cela n'est pas le cas, on en déduit que l'hypothèse
exigée par le théorème de Fubini ne tient pas~: la fonction \[
(x, y) \in [0,1] \times [0,1] \mapsto \frac{x^2 - y^2}{(x^2+y^2)^2}
\] n'est pas intégrable.

\runinhead{Homothétie}\phantomsection\addcontentsline{toc}{subsubsection}{Homothétie}\label{II-answer-h}
Si l'on pose \(D_1 = \mathbb{R}^n\), \(D_2 = \mathbb{R}^n\),
l'application \(h:D_1 \to D_2\) définie par \(h(x) = \alpha x\) est un
\(C^1\)-difféomorphisme. De plus, on a \[
|\det J_{h}(x)| = \left|\det \left[
  \begin{array}{cccc}
  \alpha & 0  & \cdots & 0 \\
  0 & \alpha & \dots & 0 \\
  \vdots & \vdots & \vdots & \vdots \\
  0 & 0 & \dots & \alpha
  \end{array} \right] \right|  = |\alpha^n| = \alpha^n.
\] Par conséquent, le
\protect\hyperref[II-theorem-changement-de-variables]{théorème de
changement de
variables (p.~\pageref*{II-theorem-changement-de-variables})} fournit \[
\int_{\mathbb{R}^n} f(y) \, dy = \int_{\mathbb{R}^n} f(\alpha x) \alpha^n \, dx,
\] soit \[
\int_{\mathbb{R}^n} f(\alpha x) \, dx = \frac{1}{\alpha^n} \int_{\mathbb{R}^n} f(y) \, dy.
\]

\runinhead{Volume et
translation}\phantomsection\addcontentsline{toc}{subsubsection}{Volume
et translation}\label{II-answer-vr}
L'ensemble \(A\) mesurable et de volume fini a une fonction
caractéristique \(1_A\) intégrable et \[
\lambda(A) = \int 1_A(x) \, dx.
\] Soit \(h(x) = x - u\) ou \(u \in \mathbb{R}^3\). La fonction \(h\)
est une bijection de \(\mathbb{R}^3\) sur lui-même, continûment
différentiable ainsi que son inverse, \(h^{-1}(x) = x + u\) et
\(J_h(x) = I\), donc \(\det J_h(x) = 1\). Par
\protect\hyperref[II-theorem-changement-de-variables]{le théorème de
changement de
variables (p.~\pageref*{II-theorem-changement-de-variables})}, la
fonction \(1_{h^{-1}(A)} = 1_{A} \circ h\) est donc intégrable sur
\(\mathbb{R}^3\) et \[
\int 1_{h^{-1}(A)}\, dx = \int_{\mathbb{R}^3} (1_A \circ h) |\det J_h(x)| \, dx = \int_{\mathbb{R}^3} 1_A(x) \, dx = \lambda(A).
\] L'ensemble translaté \(A + u = h^{-1}(A)\) est donc mesurable, de
volume fini égal au volume de \(A\).

\runinhead{Coordonnées
polaires}\phantomsection\addcontentsline{toc}{subsubsection}{Coordonnées
polaires}\label{II-answer-cp}
Les ensembles \(P\) et \(C\) sont ouverts et la fonction \(h\) -- qui
permet de passer des coordonnées polaires aux coordonnées cartésiennes
-- est une bijection de \(P\) dans \(C\). Elle est continûment
différentiable\,; sa matrice jacobienne en \((r, \theta)\) vaut \[
J_{h}(r, \theta) = 
\left[ 
\begin{array}{cr}
\cos \theta & -r \sin \theta \\
\sin \theta & r \cos \theta
\end{array}
\right],
\] et son déterminant satisfait \[
\det J_{h}(r, \theta) = (\cos \theta)(r \cos \theta) - (\sin \theta)(-r\sin \theta)
= r > 0.
\] Le jacobien est donc inversible et \(h^{-1}\) est continûment
différentiable par le théorème d'inversion locale. On peut donc
appliquer \protect\hyperref[II-theorem-changement-de-variables]{le
théorème de changement de
variables (p.~\pageref*{II-theorem-changement-de-variables})} à la
fonction \(f\), ce qui fournit \begin{align*}
\int_C f(x, y) \, d(x, y) &= \int_P f(h(r, \theta)) |\det J_h(r,\theta)| d(r,\theta) \\
&= \int_P g(r, \theta)  r \, d(r,\theta).
\end{align*}

\runinhead{Absence du déterminant
jacobien}\phantomsection\addcontentsline{toc}{subsubsection}{Absence du
déterminant jacobien}\label{II-answer-adj}
La façon la plus rapide de procéder consiste à considérer que
\(f \circ h\) joue le rôle de \(f\) dans
\protect\hyperref[II-theorem-changement-de-variables]{le théorème de
changement de
variables (p.~\pageref*{II-theorem-changement-de-variables})}, que \(h\)
dans notre énoncé désigne \(h^{-1}\) dans ce théorème et que les rôles
de \(D_1\) et \(D_2\) sont intervertis. Une fois que l'on a permuté ces
notations, on réalise que l'intégrale que l'on souhaite calculer est le
membre de gauche de l'équation du théorème de changement de variables,
dont toutes les hypothèses sont par ailleurs satisfaites. Par conséquent
on a \[
\int_{D_1} f (h(x)) \, dx = \int_{D_2} (f \circ h)(h^{-1}(y)) |\det J_{h^{-1}}(y)| \, dy.
\] On peut simplifier \((f \circ h)(h^{-1}(y))\) en \(f(y)\) et
éventuellement exprimer le jacobien de \(h^{-1}\) en fonction du
jacobien de \(h\)~: \(J_{h^{-1}}(y) = [J_{h}(h^{-1}(y))]^{-1}\)\,; on
obtient donc \[
\int_{D_1} f (h(x)) \, dx = 
\int_{D_2} f(y) |\det J_{h^{-1}}(y)| \, dy = 
\int_{D_2} f(y) \frac{1}{|\det J_{h}(h^{-1}(y))|} \, dy.
\]

\runinhead{Lancer de
fléchette}\phantomsection\addcontentsline{toc}{subsubsection}{Lancer de
fléchette}\label{II-answer-exo-flechette}
Les coordonnées cartésiennes de
\(M \in D = \{(x,y) \in \mathbb{R}^2 , x^2+y^2 \leq 1\}\) constituent un
couple de variables aléatoires de densité
\[f_{(X,Y)}(x,y) = \frac{1}{\pi}1_{D} (x,y)\] uniforme sur le disque,
par hypothèse. L'abscisse de \(X\) est distribuée selon la densité
marginale
\[f_X(x) = \int f_{(X,Y)}(x,y) dy = \frac{2}{\pi} (1-x^2)^{1/2} 1_{[-1,1]}(x).\]
La loi de \(Y\) a la même densité.

\runinhead{Lancer de fléchette
(suite)}\phantomsection\addcontentsline{toc}{subsubsection}{Lancer de
fléchette (suite)}\label{II-answer-exo-flechettebis}
La densité des coordonnées cartésiennes ne s'exprime pas comme le
produit des densités marginales\,; elles ne sont pas indépendantes.

\runinhead{Démonstration}\phantomsection\addcontentsline{toc}{subsubsection}{Démonstration}\label{II-answer-demo-cov-indep}
Soit \(X\) et \(Y\) deux v.a. aléatoires indépendantes et de carré
intégrable. On a d'après \protect\hyperref[II-indep_fct]{la
proposition (p.~\pageref*{II-indep_fct})}
\[\mathbb{E}((X-\mathbb{E}(X))(Y-\mathbb{E}(Y))) = \mathbb{E}(X-\mathbb{E}(X))\mathbb{E}(Y-\mathbb{E}(Y)) = 0 = \rho(X,Y)\]

\runinhead{Lancer de fléchette
(fin)}\phantomsection\addcontentsline{toc}{subsubsection}{Lancer de
fléchette (fin)}\label{II-answer-exo-flechetteter}
En coordonnées polaires, le domaine \(D\) est décrit par
\(D = \{(r,\theta) \in [0,1] \times [0,2\pi]\}\). La densité de
\((R,\theta)\) s'écrit ainsi~:
\[f_{(R,\theta)}(r,\theta) = \frac{r}{\pi}1_{D} (r\cos(\theta),r\sin(\theta)) = 2r 1_{[0,1]}(r) \frac{1}{2\pi}1_{[0,2\pi]}(r).\]
Ainsi \(R\) et \(\Theta\) sont indépendantes de densités respectives
\(2r 1_{[0,1]}(r)\) et \(\frac{1}{2\pi}1_{[0,2\pi]}(\theta)\).

\runinhead{Loi
bêta}\phantomsection\addcontentsline{toc}{subsubsection}{Loi
bêta}\label{II-answer-exo-loibeta}
On note d'abord que la dimension de \(Y\) est plus petite que celle de
\(X\). On va donc compléter \(Y\) en prenant par exemple \(Y' = (Y,Z)\),
avec \(Z=U+V\), ce qui correspond à
\(g(u,v) = \left(\frac{u}{u+v},u+v\right)\). Cette application est
bijective de \(A = \left]0,+\infty\right[^2\) dans
\(B = \left]0,1\right[ \times \left]0,+\infty\right[\), d'inverse
\(g^{-1}(y,z) = (yz,z(1-y))\), qui a pour jacobien \(z\).

Comme
\(f_X(u,v) = \frac{\theta^{\alpha+\beta}}{\Gamma(\alpha)\Gamma(\beta)}u^{\alpha-1}v^{\beta-1}e^{-\theta(u+v)}1_A(u,v)\),
on a
\[f_{Y'}(y,z) = \frac{\theta^{\alpha+\beta}}{\Gamma(\alpha)\Gamma(\beta)}z^{\alpha+\beta -1}y^{\alpha-1}(1-y)^{\beta-1}e^{-\theta z}1_B(y,z).\]
On obtient alors la densité de Y en intégrant \(f_{Y'}(y,z)\) par
rapport à \(z\in \left]0,+\infty\right[\)~: \begin{align*}
f_Y(y) &= \int_0^{+\infty} f_{Y'}(y,z) dz\\
       &= \frac{\theta^{\alpha+\beta}}{\Gamma(\alpha)\Gamma(\beta)}y^{\alpha-1}(1-y)^{\beta-1}1_{\left]0,1\right[}(y)\int_0^{+\infty}z^{\alpha+\beta -1}e^{-\theta(z)}dz\\
       &= \frac{\Gamma(\alpha+\beta)}{\Gamma(\alpha)\Gamma(\beta)}y^{\alpha-1}(1-y)^{\beta-1}1_{\left]0,1\right[}(y),
\end{align*} où on a utilisé la définition de la fonction gamma et le
changement de variable linéaire \(z \mapsto \theta z\). On appelle loi
bêta de paramètres \(\alpha\) et \(\beta\) la loi admettant cette
densité. Admettant une grande variété de formes, elle permet de
modéliser de nombreuses distributions à support fini.

La loi bêta apparaît naturellement dans une expérience d'urnes, donnée
par George Pólya dans un article de 1930,
\href{http://www.numdam.org/article/AIHP_1930__1_2_117_0.pdf}{\emph{Sur
quelques points de la théorie des probabilités}}. Il décrit l'expérience
suivante~: on se donne une urne contenant initialement \(r\) boules
rouges et \(b\) boules bleues, on tire une boule dans l'urne, puis on la
remet dans l'urne avec une deuxième boule de même couleur. Alors la
proportion de boules rouges tend vers une variable aléatoire de loi
Beta\((r,b)\), et, inversement, la proportion de boules bleues tend vers
une variable aléatoire de loi Beta\((b,r)\).

Nous obtenons aussi facilement la densité de \(Z\). En effet, on a
\(f_{Y'}(y,z) = f_Y(y)f_Z(z)\) (\(Y\) et \(Z\) sont donc indépendantes),
où
\[f_Z(z) = \frac{\theta^{\alpha+\beta}}{\Gamma(\alpha+\beta)}z^{\alpha+\beta -1}e^{-\theta z}1_{\left]0,+\infty\right[}\]
On a ainsi démontré que si \(U\) et \(V\) sont deux variables aléatoires
indépendantes de lois respectives \(\Gamma(\alpha,\theta)\) et
\(\Gamma(\beta,\theta)\), alors \(U+V\) suit la loi
\(\Gamma(\alpha+\beta,\theta)\) et est indépendante de \(\frac{U}{U+V}\)
qui suit une loi bêta de paramètres \((\alpha,\beta)\).

\runinhead{Divergence de Kullback-Leibler entre deux
gaussiennes}\phantomsection\addcontentsline{toc}{subsubsection}{Divergence
de Kullback-Leibler entre deux gaussiennes}\label{II-answer-exoklgauss}
On a
\[f(x) = \frac{1}{\sqrt{2\pi}\sigma_1} e^{-\frac{(x-\mu_1)^2}{2\sigma_1^2}}\]
et
\[g(x) = \frac{1}{\sqrt{2\pi}\sigma_2} e^{-\frac{(x-\mu_2)^2}{2\sigma_2^2}}\]

On en déduit que \begin{align*}
D_{KL} (f || g) &= \int_{-\infty}^{+\infty} f(x) \log \left(\frac{f(x)}{g(x)}\right) dx \\
                &= \int_{-\infty}^{+\infty} f(x) \left[ \log \left(\frac{\sigma_2}{\sigma_1}\right) + \frac{(x-\mu_2)^2}{2\sigma_2^2} - \frac{(x-\mu_1)^2}{2\sigma_1^2} \right] dx \\
                &= \log \left(\frac{\sigma_2}{\sigma_1}\right) + \frac{1}{2\sigma_2^2} \int_{-\infty}^{+\infty} f(x) (x-\mu_2)^2 dx - \frac{1}{2\sigma_1^2} \int_{-\infty}^{+\infty} f(x) (x-\mu_1)^2 dx
\end{align*} On a~: \begin{align*}
\int_{-\infty}^{+\infty} f(x) (x-\mu_2)^2 dx &= \int_{-\infty}^{+\infty} f(x) ((x-\mu_1)+\mu_1-\mu_2)^2 dx \\
&= \int_{-\infty}^{+\infty} f(x) (x-\mu_1)^2 dx + (\mu_1 - \mu_2)^2 \\
&= \sigma_1^2 + (\mu_1 - \mu_2)^2 
\end{align*} On en déduit la formule annoncée~:
\[D_{KL} (f || g) = \log \left(\frac{\sigma_2}{\sigma_1}\right) + \frac{\sigma_1^2 + (\mu_1 - \mu_2)^2}{2\sigma_2^2} - \frac{1}{2}\]

\section{Durée de vie}\label{II-duruxe9e-de-vie-1}

\runinhead{Question
1}\phantomsection\addcontentsline{toc}{subsubsection}{Question
1}\label{II-answer-dureevie1}
En dérivant, on obtient que la densité de \(T\) est
\(f(t) = t e^{-t^2/2}1_{t>0}\).

L'espérance vaut~:

\begin{align*}
\mathbb{E}(T) &= \int_0^{+\infty} t^2 e^{-t^2/2} dt\\
        &= \left[-t e^{-t^2/2} \right]_0^{+\infty} + \int_0^{+\infty} e^{-t^2/2} dt \text{ par intégration par parties}\\
        &= \frac{1}{2} \int_{-\infty}^{+\infty} e^{-t^2/2} dt \\
        &= \frac{\sqrt{2\pi}}{2}
\end{align*}

\runinhead{Question
2}\phantomsection\addcontentsline{toc}{subsubsection}{Question
2}\label{II-answer-dureevie2}
La probabilité s'écrit~:

\[\mathbb{P}(T\geq 3 |T \geq 1) = \frac{\mathbb{P}(T \geq 3,T \geq 1)}{\mathbb{P}(T \geq 1)} = \frac{\mathbb{P}(T \geq 3)}{\mathbb{P}(T \geq 1)} = \frac{e^{-9/2}}{e^{-1/2}} = e^{-4}\]

\runinhead{Question
3}\phantomsection\addcontentsline{toc}{subsubsection}{Question
3}\label{II-answer-dureevie3}
Les variables aléatoires \(X_1 , \ldots , X_{10}\) sont indépendantes et
suivent une loi de Bernoulli de paramètre~:
\[ \mathbb{P}(T \leq 1) = F (1) = 1 - e^{- 1/2}\]

On en déduit que la loi de N est la loi binomiale de paramètre
\((10, 1 - e^{- 1/2})\).

\runinhead{Question
4}\phantomsection\addcontentsline{toc}{subsubsection}{Question
4}\label{II-answer-dureevie4}
La probabilité que l'équipement en série soit défaillant avant 1 an
vaut~:
\[ P(N \geq 1) = 1 - \mathbb{P}(N = 0) = 1 - e^{-10/2} \approx 9.9 ~10^{-1} = 99\%\]

\runinhead{Question
5}\phantomsection\addcontentsline{toc}{subsubsection}{Question
5}\label{II-answer-dureevie5}
La probabilité que l'équipement en parallèle soit défaillant avant 1 an
vaut~:
\[ \mathbb{P}(N = 10) = (1 - e^{-1/2})^10 \approx 8.9 ~10 ^{-5} \]

\section{Crues centennales de la
Seine}\label{II-crues-centennales-de-la-seine-1}

\runinhead{Question
1}\phantomsection\addcontentsline{toc}{subsubsection}{Question
1}\label{II-answer-crue1}
Soit \(x\in \mathbb{R}\) \begin{align*}
F_Y(x) = F(Y\leq x) &= F(X_1\leq x,\ldots,X_n\leq x)\\
           &= \prod_{i=1}^n F(X_i\leq x) \\
           &= F(x)^n
\end{align*} puisque les \(X_i\) sont indépendantes et de même loi

\runinhead{Question
2}\phantomsection\addcontentsline{toc}{subsubsection}{Question
2}\label{II-answer-crue2}
Soit \(x\in \mathbb{R}\) \begin{align*}
f_Y(x) &= \frac{d}{d x} F_Y(x)\\
       &= \frac{d}{d x} F(x)^n\\
       &= nf(x)F(x)^{n-1}
\end{align*}

\runinhead{Question
3}\phantomsection\addcontentsline{toc}{subsubsection}{Question
3}\label{II-answer-crue3}
On a alors \(f(x) = 1_{[0,1]}(x)\) et
\(F(x) = x 1_{[0,1]}(x) + 1_{[1,+\infty]}\). On en déduit
\[ f_Y(x) = n x^{n-1}1_{[0,1]} \]

\runinhead{Question
4}\phantomsection\addcontentsline{toc}{subsubsection}{Question
4}\label{II-answer-crue4}
On cherche à déterminer \(x\) tel que
\(\mathbb{P}(Y > x) \leq \frac{1}{1000}\). Or,
\[ \mathbb{P}(Y > x) = 1-\mathbb{P}(Y \leq x) = 1 - F_Y(x) = 1 - F(x)^n = 1-x^n \]
Ainsi
\[ 1-x^n \leq \frac{1}{1000} \Leftrightarrow (1-\frac{1}{1000})^{1/n} \leq x \]

\runinhead{Question
5}\phantomsection\addcontentsline{toc}{subsubsection}{Question
5}\label{II-answer-crue5}
On rappelle que la médiane est le réel \(m\) tel que
\(\frac{1}{2} = \mathbb{P}(Y \leq m) = F_Y(m)\).

On cherche donc \(m\) tel que \(m^n = \frac{1}{2}\), soit
\(m = \frac{1}{2^{1/n}}\).

On voit que cette valeur est plus élevée que la médiane des \(X_i\) qui
vaut \(\frac{1}{2}\). Elle tend même vers 1 lorsque \(n\) tend vers
l'infini.

\section{Indépendance et vecteurs
gaussiens}\label{II-induxe9pendance-et-vecteurs-gaussiens-1}

\runinhead{Question
1}\phantomsection\addcontentsline{toc}{subsubsection}{Question
1}\label{II-answer-covindepgauss-joint}
D'après le cours, en interprétant les éléments de \(\mathbb{R}^2\) comme
des vecteurs colonnes, pour tout \(z =(x,y)\in\mathbb{R}^2\) on a \[
f_Z(z) = \dfrac{1}{2\pi\,\sqrt{\det(C)}}\,\exp\left\{ -\dfrac{1}{2}\,\left(z - m\right)^t\,C^{-1}\,\left(z - m\right) \right\}.
\] Or par construction \(C\) peut s'écrire
\[C = \left( \begin{array}{cc} \sigma_X^2 & c\\ c & \sigma_Y^2 \end{array}\right), \]
où \(\sigma_X^2 := \mathbb{V}\left(X\right)\),
\(\sigma_Y^2 := \mathbb{V}\left(Y\right)\) et
\(c = \text{Cov}\left(X,Y\right)\). Comme \(C\) est définie positive
elle est inversible, avec \(\det(C) = \sigma_X^2\,\sigma_Y^2 - c^2 > 0\)
et
\[C^{-1} = \dfrac{1}{\sigma_X^2\,\sigma_Y^2 - c^2}\,\left( \begin{array}{cc} \sigma_Y^2 & -c\\ -c & \sigma_X^2 \end{array}\right).\]
En outre, \(m = \left(m_X, m_Y \right)\) avec \(m_X := \mathbb{E}(X)\)
et \(m_Y := \mathbb{E}(Y)\). On en déduit que \begin{align*}
f_Z(x,y) &= \dfrac{1}{2\pi\,\sqrt{\sigma_X^2\,\sigma_Y^2 - c^2}}\,\exp\biggl\{ \dfrac{-1}{2\,\left(\sigma_X^2\,\sigma_Y^2 - c^2 \right)}\times\\
&\ \ \ \left( \sigma_Y^2\,(x-m_X)^2 - 2c\,(x-m_X)\,(y-m_Y) + \sigma_X^2\,(y-m_Y)^2 \right) \biggr\}\\
&= \dfrac{1}{2\pi\,\sigma_X\,\sigma_Y\,\sqrt{1-\rho^2}}\,\,\exp\biggl\{ \dfrac{-1}{2\,\left( 1 - \rho^2 \right)}\times\\
&\ \ \ \left( \dfrac{(x-m_X)^2}{\sigma_X^2} - 2\rho\,\dfrac{(x-m_X)}{\sigma_X}\,\dfrac{(y-m_Y)}{\sigma_Y} + \dfrac{(y-m_Y)^2}{\sigma_Y^2} \right) \biggr\}.
\end{align*}

\runinhead{Question
2}\phantomsection\addcontentsline{toc}{subsubsection}{Question
2}\label{II-answer-covindepgauss-marg}
Ici, pour \(x,y \in \mathbb{R}^2\), on remarque que~: \begin{align*}
f_Z(x,y) &= \dfrac{1}{2\pi\,\sigma_X\,\sigma_Y\,\sqrt{1-\rho^2}}\,\,\exp\biggl\{ \dfrac{-1}{2\,\left( 1 - \rho^2 \right)}\times\\
&\ \ \ \left( \left(\dfrac{y-m_Y}{\sigma_Y} - \rho\,\dfrac{x-m_X}{\sigma_X}\right)^2 + (1 - \rho^2)\,\dfrac{(x-m_X)^2}{\sigma_X^2} \right) \biggr\}\\
& = \dfrac{1}{\sqrt{2\pi}\,\sigma_X}\,\exp\left\{-\dfrac{(x-m_X)^2}{2\,\sigma_X^2} \right\}\times\\
&\ \ \ \dfrac{1}{\sqrt{2\pi}\,\sigma_Y\,\sqrt{1-\rho^2}}\,\exp\left\{-\dfrac{1}{2\,(1-\rho^2)}\,\left(\dfrac{y-m_Y}{\sigma_Y} - \rho\,\dfrac{x-m_X}{\sigma_X}\right)^2 \right\}.
\end{align*} Ainsi, avec le changement de variable
\(u = \dfrac{y-m_Y}{\sigma_Y}\) on obtient \begin{align*}
f_X(x) &= \int_{-\infty}^{+\infty} f_Z(x,y)\,dy\\
&= \dfrac{1}{\sqrt{2\pi}\,\sigma_X}\,\exp\left\{-\dfrac{(x-m_X)^2}{2\,\sigma_X^2} \right\}\times\\
&\ \ \ \int_{-\infty}^{+\infty} \dfrac{1}{\sqrt{2\pi}\,\sqrt{1-\rho^2}}\,\exp\left\{-\dfrac{1}{2\,(1-\rho^2)}\,\left(u - \rho\,\dfrac{x-m_X}{\sigma_X}\right)^2 \right\}\,du,
\end{align*} et on reconnaît dans cette dernière intégrale la densité
d'une loi Normale d'espérance \(\rho\,\dfrac{x-m_X}{\sigma_X}\) et de
variance \(1-\rho^2\). Cette intégrale vaut donc \(1\) et on conclut que
\[f_X(x) = \dfrac{1}{\sqrt{2\pi}\,\sigma_X}\,\exp\left\{-\dfrac{(x-m_X)^2}{2\,\sigma_X^2} \right\}, \]
qui correspond à la densité d'une loi Normale d'espérance \(m_X\) et de
variance \(\sigma_X^2\).

En procédant de manière symétrique, on obtient de même que \(Y\) suit
une loi Normale d'espérance \(m_Y\) et de variance \(\sigma_Y^2\).

\runinhead{Question
3}\phantomsection\addcontentsline{toc}{subsubsection}{Question
3}\label{II-answer-covindepgauss-indep}
Le premier sens est évident~: si \(X\) et \(Y\) sont indépendantes,
alors nous avons vu dans le cours que
\(\text{Cov}\left(X,Y\right) = 0\), et ce que l'on soit gaussien ou non.

Supposons maintenant que \(\text{Cov}\left(X,Y\right) = 0\), et montrons
l'indépendance entre \(X\) et \(Y\). En reprenant la formule de la
question 1 et en remplaçant \(\rho\) par \(0\), pour tout
\((x,y) \in \mathbb{R}^2\) on a: \begin{align*}
f_Z(x,y) &= \dfrac{1}{2\,\pi\,\sigma_X\,\sigma_Y}\,\exp\left\{ -\dfrac{1}{2}\,\left( \dfrac{(x-\mu_X)^2}{\sigma_X^2} + \dfrac{(y-\mu_Y)^2}{\sigma_Y^2} \right) \right\}\\
&= \dfrac{1}{\sqrt{2\pi}\,\sigma_X}\,\exp\left\{ -\dfrac{(x-m_X)^2}{2\,\sigma_X^2} \right\} \times \dfrac{1}{\sqrt{2\pi}\,\sigma_Y}\,\exp\left\{ -\dfrac{(y-m_Y)^2}{2\,\sigma_Y^2} \right\}\\
&= f_X(x)\times f_Y(y)
\end{align*} d'après la question 2. On en conclut que \(X\) et \(Y\)
sont bien indépendantes.

\section{Symétrie de la
gaussienne}\label{II-symuxe9trie-de-la-gaussienne-1}

\runinhead{Question
1}\phantomsection\addcontentsline{toc}{subsubsection}{Question
1}\label{II-answer-gausssym-loi}
Nous avons admis dans le cours que la fonction de répartition \(F_c\) de
\(X_c\) caractérisait sa loi. Nous allons donc l'expliciter. Soit
\(x \in \mathbb{R}\). Par définition de la fonction de répartition d'une
variable aléatoire et d'après la formule des probabilités totales, on a
\begin{align*}
F_c(x) &= \mathbb{P}\left(X_c \leq x\right) = \mathbb{P}\left( X_c \leq x, |X| > c \right) + \mathbb{P}\left( X_c \leq x, |X| \leq c\right)\\
&= \mathbb{P}\left( X\leq x, |X| > c \right) + \mathbb{P}\left(-X \leq x, |X| \leq c\right)\\
&= \mathbb{P}\left( X\leq x, X > c \right) + \mathbb{P}\left(X \leq x, X < -c \right) + \mathbb{P}\left( X \geq -x, -c \leq X \leq c\right).
\end{align*} Or
\[\mathbb{P}\left( X\leq x, X > c \right) = \left|\begin{array}{ll} \mathbb{P}\left(c < X \leq x \right) = F(x) - F(c) & \text{si } x > c,\\ 0 & \text{si } x\leq c,\end{array}\right.\]
\[ \mathbb{P}\left(X \leq x, X < -c \right) = \mathbb{P}\left( X \leq \min\{-c,x\} \right) = F\left(\min\{-c,x\}\right), \]
\[ \mathbb{P}\left( X \geq -x, -c \leq X \leq c\right) = \left|\begin{array}{ll} \mathbb{P}\left(-c \leq X \leq c \right) = F(c) - F(-c) & \text{si } x > c, \\ \mathbb{P}\left(-x \leq X \leq c \right) = F(c) - F(-x) & \text{si } -c \leq x \leq c,\\ 0 & \text{si } x<-c. \end{array}\right.\]
Ainsi,

\begin{itemize}
\item
  si \(x > c\) alors \(\min\{-c,x\} = -c\) et
  \[F_c(x) = F(x) - F(c) + F\left(\min\{-c,x\}\right) + F(c) - F(-c) = F(x),\]
\item
  si \(-c \leq x \leq c\) alors \(\min\{-c,x\} = -c\) et
  \[F_c(x) = F\left(\min\{-c,x\}\right) + F(c) - F(-x) = 1 - F(c) + F(c) - 1 + F(x) = F(x)\]
  car la densité de la loi Normale centrée réduite est paire, ce qui
  implique que \(F(-x) = 1 - F(x)\) pour tout \(x\in\mathbb{R}\),
\item
  si \(x < -c\) alors \(\min\{-c,x\} = x\) et
  \[F_c(x) =  F\left(\min\{-c,x\}\right) = F(x).\]
\end{itemize}

On en conclut que \(F_c = F\)~: \(X_c\) a la même loi que \(X\). Ce
résultat est dû à la symétrie de la densité de la loi Normale centrée
réduite. On remarque d'ailleurs que le résultat resterait inchangé si
l'on prenait n'importe quelle autre loi de densité paire\,!

\runinhead{Question
2}\phantomsection\addcontentsline{toc}{subsubsection}{Question
2}\label{II-answer-gausssym-cov}
Tout d'abord, par définition,
\[\text{Cov}\left(X,X_c\right) = \mathbb{E}\left( \left(X - \mathbb{E}\left(X\right)\right)\,\left(X_c - \mathbb{E}\left(X_c\right) \right) \right) = \mathbb{E}\left(X\,X_c \right) \]
puisque \(X\) et \(X_c\) sont gaussiennes centrées (d'espérance nulle)
réduites. Or on remarque que l'on peut écrire
\(X_c = X\,1_{\mathbb{R}\backslash[-c,c]}(X) - X\,\,1_{[-c,c]}(X)\). On
déduit par linéarité de l'espérance que \begin{align*}
\mathbb{E}\left(X\,X_c\right) &= \mathbb{E}\left(X^2\,1_{\mathbb{R}\backslash[-c,c]}(X) - X^2\,1_{[-c,c]}(X)\right)\\
&= \mathbb{E}\left(X^2\,\left(1_{\mathbb{R}\backslash[-c,c]}(X) + 1_{[-c,c]}(X) - 1_{[-c,c]}(X)\right) - X^2\,1_{[-c,c]}(X)\right)\\
&= \mathbb{E}\left(X^2\right) - 2\,\mathbb{E}\left(X^2\,1_{[-c,c]}(X)\right).
\end{align*} Etant donné que \(X\) est de loi Normale centrée réduite,
on a directement que \(\mathbb{E}\left(X^2\right)=1\); il suffit donc de
calculer la dernière espérance pour obtenir le résultat. Comme
\(x\mapsto x^2\,1_{[-c,c]}(x)\) est bornée et continue sur
\(\mathbb{R}\backslash\{-c,c\}\) (donc continue presque partout), elle
est intégrable et
\[\mathbb{E}\left(X^2\,1_{[-c,c]}(X)\right) = \int_{-\infty}^{+\infty} x^2\,1_{[-c,c]}\,f(x)\,dx = \int_{-c}^c x^2\,f(x)\,dx. \]
Or on pourra remarquer que \(f^\prime(x) = -x\,f(x)\) puis que
\(f^{\prime\prime}(x) = (x^2-1)\,f(x)\). Puisque \(f\), \(f^\prime\) et
\(f^{\prime\prime}\) sont continues et bornées (voir l'exercice
\emph{Densité et fonction de répartition d'une loi Normale} du chapitre
Probabilités I) elles sont intégrables sur des segments. On en déduit
que \begin{align*}
\mathbb{E}\left(X^2\,1_{[-c,c]}(X)\right) &= \int_{-c}^c (x^2-1)\,f(x)\,dx + \int_{-c}^c f(x)\,dx\\
&= \left[-x\,f(x) \right]_{-c}^c + F(c) - F(-c)\\
&= -c\,\left(f(c) + f(-c)\right) + F(c) - F(-c)\\
&= 2\,\left(F(c) - c\,f(c)\right) - 1
\end{align*} par parité de \(f\). On en conclut que
\(\text{Cov}\left(X,X_c\right) = 3 + 4\,\left(c\,f(c) - F(c)\right)\).

La fonction
\(g : c\in\mathbb{R}_+^\ast \mapsto 3+4\,\left(c\,f(c) - F(c)\right)\)
est continue, infiniment dérivable et pour tout \(c > 0\) on a
\(g^\prime(c) = -4\,c^2\,f(c) < 0\). Par conséquent, \(g\) est
strictement décroissante, avec \(\lim\limits_{c\to0^+} g(c) = 1\) et
\(\lim\limits_{c\to+\infty} g(c) = -1\). Le théorème des valeurs
intermédiaires nous permet de conclure qu'il existe bien un unique
\(c_0 > 0\) tel que \(g(c_0) = 0\).

\runinhead{Question
3}\phantomsection\addcontentsline{toc}{subsubsection}{Question
3}\label{II-answer-gausssym-indep}
Les variables aléatoires \(X\) et \(X_{c_0}\) sont indépendantes ssi
pour tous ensembles \(A,B \in \mathcal{E}_\mathbb{R}\) on a
\[\mathbb{P}\left(X \in A, X_{c_0} \in B \right) = \mathbb{P}\left(X \in A\right)\,\mathbb{P}\left(X_{c_0} \in B \right).\]

Or si l'on prend \(A = B = ]-\infty, x]\) où \(x < -c_0\), alors
\begin{align*}
\mathbb{P}\left( X \leq x, X_{c_0} \leq x \right) &= \mathbb{P}\left( X \leq x, X < -c_0 \right) + \mathbb{P}\left( X \leq x, X > c_0 \right)\\
&\ \ \ + \mathbb{P}\left( X\leq x, X \geq -x, -c_0 \leq X \leq c_0 \right)\\
&= \mathbb{P}\left( X \leq x \right) +0 + 0 = F(x)\\
& \neq F(x)\,F(x) = \mathbb{P}\left( X \leq x \right)\,\mathbb{P}\left(X_{c_0} \leq x \right).
\end{align*}

On en conclut que bien que \(\text{Cov}\left(X,X_{c_0}\right) = 0\),
\(X\) et \(X_{c_0}\) ne sont pas indépendantes.

\runinhead{Question
4}\phantomsection\addcontentsline{toc}{subsubsection}{Question
4}\label{II-answer-gausssym-vecgauss}
Supposons que \((X,X_{c_0})\) soit un vecteur gaussien. Alors, puisque
\(\text{Cov}\left(X,X_{c_0}\right) = 0\), \(X\) et \(X_{c_0}\) sont
nécessairement indépendantes. Or nous venons de voir que ce n'était pas
le cas; nous avons donc contradiction. On en conclut que \((X,X_{c_0})\)
n'est pas gaussien.

\section{Loi du \(\chi^2\)}\label{II-loi-du-chi2-1}

\runinhead{A 1 degré de
liberté}\phantomsection\addcontentsline{toc}{subsubsection}{A 1 degré de
liberté}\label{II-answer-khi2-1dl}

\begin{enumerate}
\def\labelenumi{\arabic{enumi}.}
\tightlist
\item
  Soit \(h : \mathbb{R}\to \mathbb{R}\) une fonction mesurable telle que
  \(h\times f\) est intégrable. Alors en opérant le changement de
  variable \(u = x^2\) on a \begin{align*}
  \mathbb{E}\left( h(Y) \right) &= \mathbb{E}\left(h(X^2)\right) = \int_{\mathbb{R}} h(x^2)\,f(x)\,dx\\
  & = \int_{-\infty}^{0} h(x^2)\,f(x)\,dx + \int_{0}^{+\infty} h(x^2)\,f(x)\,dx\\
  &= \int_{0}^{+\infty} h(u)\,f\left(-\sqrt{u}\right)\,\dfrac{1}{2\,\sqrt{u}} \,du + \int_{0}^{+\infty} h(u)\,f\left(\sqrt{u}\right)\,\dfrac{1}{2\,\sqrt{u}} \,du.
  \end{align*} Or \(f\) est paire, donc pour tout \(u \in \mathbb{R}^+\)
  on a \(f\left(-\sqrt{u}\right) = f\left(\sqrt{u}\right)\) et
  \begin{align*}
  \mathbb{E}\left( h(Y) \right) &= \int_{\mathbb{R}_+} h(u)\,f\left(\sqrt{u}\right)\,\dfrac{1}{\sqrt{u}} \,du.
  \end{align*} Par propriété \(Y\) possède donc une densité \(f_Y\)
  valant pour tout \(x\in\mathbb{R}\)
  \[f_Y(x) = \left|\begin{array}{ll} f\left(\sqrt{x}\right)\,\dfrac{1}{\sqrt{x}} = \dfrac{x^{-\frac{1}{2}}}{2^{\frac{1}{2}}\,\Gamma\left(\frac{1}{2}\right)}\,\exp\left\{-\dfrac{x}{2}\right\} & \text{si } x > 0,\\ 0 & \text{si } x \leq 0. \end{array}\right. \]
\end{enumerate}

\textbf{Remarque.} La fonction \(f_Y\) est bien une densité~: elle est
continue par morceaux donc mesurable, puis positive, et son intégrale
sur \(\mathbb{R}\) vaut bien \(1\)~: \begin{align*}
\int_{\mathbb{R}} f_Y(x)\,dx &= \dfrac{1}{2\,\Gamma\left(\frac{1}{2}\right)}\,\int_{0}^{+\infty} \left(\dfrac{x}{2}\right)^{\frac{1}{2}-1}\,\exp\left\{-\dfrac{x}{2}\right\}\,dx \\
&= \dfrac{1}{\Gamma\left(\frac{1}{2}\right)}\,\int_{0}^{+\infty} u^{\frac{1}{2}-1}\,e^{-u}\,du = \dfrac{\Gamma\left(\frac{1}{2}\right)}{\Gamma\left(\frac{1}{2}\right)} = 1.
\end{align*}

\runinhead{A n degrés de
liberté}\phantomsection\addcontentsline{toc}{subsubsection}{A n degrés
de liberté}\label{II-answer-khi2-ndl}

\begin{enumerate}
\def\labelenumi{\arabic{enumi}.}
\setcounter{enumi}{1}
\tightlist
\item
  On note \((\mathcal{P}_n)\) la propriété à démonter au rang \(n\).
\end{enumerate}

\textbf{Initialisation.} Nous avons vu à la question précédente que
\((\mathcal{P}_1)\) est vraie.

\textbf{Hérédité.} Soit maintenant \(n \geq 2\) et supposons
\((\mathcal{P}_{n-1})\) vraie. Alors \(Y = \sum_{i = 1}^{n-1}Y_i + Y_n\)
est la somme de deux variables aléatoires indépendantes, dont la
première suit un \(\chi_{n-1}^2\) de densité notée \(f_{n-1}\) et la
seconde un \(\chi_1^2\) de densité notée \(f_1\). Elle possède donc une
densité, égale au produit de convolution de \(f_{n-1}\) et \(f_1\)~:
pour tout \(x\in\mathbb{R}\)
\[f_Y(x) = \int_{\mathbb{R}} f_1(x-u)\,f_{n-1}(u)\,du.\] Or pour tout
\(x,u\in\mathbb{R}\) on a \(f_1(x-u) = 0\) ssi \(x \leq u\) et
\(f_{n-1}(u) = 0\) ssi \(u \leq 0\). Donc leur produit est strictement
positif ssi \(x > u > 0\), nul sinon. La densité \(f_Y\) est par
conséquent nulle sur \(\mathbb{R}^-\) et nous allons maintenant exhiber
son expression sur \(\mathbb{R}_+^\ast\). Soit \(x > 0\), alors
\begin{align*}
f_Y(x) &= \int_{0}^{x} \dfrac{(x-u)^{-\frac{1}{2}}}{\sqrt{2}\,\Gamma\left(\frac{1}{2}\right)}\,\exp\left\{-\dfrac{x-u}{2}\right\}\,\dfrac{u^{\frac{n-1}{2}-1}}{2^{\frac{n-1}{2}}\,\Gamma\left(\frac{n-1}{2}\right)}\,\exp\left\{-\dfrac{u}{2}\right\} \,du\\
&= \dfrac{\exp\left\{-\dfrac{x}{2} \right\}}{2^{\frac{n}{2}}\,\Gamma\left(\frac{1}{2}\right)\Gamma\left(\frac{n-1}{2}\right)} \,\int_0^x \left(x-u\right)^{-\frac{1}{2}}\,u^{\frac{n-3}{2}}\,du.
\end{align*} Avec le changement de variable \(v=\dfrac{u}{x}\) on
obtient \begin{align*}
f_Y(x) &= \dfrac{x^{\frac{n}{2}-1}\,\exp\left\{-\dfrac{x}{2} \right\}}{2^{\frac{n}{2}}\,\Gamma\left(\frac{1}{2}\right)\Gamma\left(\frac{n-1}{2}\right)} \,\int_0^{1} v^{\frac{n-3}{2}}\,\left(1 - v\right)^{-\frac{1}{2}} \,dv.
\end{align*} Or \(f_Y\) est une densité, son intégrale vaut donc \(1\)~:
\begin{align*}
\int_{\mathbb{R}} f_Y(x)\,dx &= \int_0^{+\infty} \dfrac{x^{\frac{n}{2}-1}\,\exp\left\{-\dfrac{x}{2} \right\}}{2^{\frac{n}{2}}\,\Gamma\left(\frac{1}{2}\right)\Gamma\left(\frac{n-1}{2}\right)}\,\left(\int_0^{1} v^{\frac{n-3}{2}}\,\left(1 - v\right)^{-\frac{1}{2}} \,dv\right)\,dx\\
&= \left(\int_0^{1} v^{\frac{n-3}{2}}\,\left(1 - v\right)^{-\frac{1}{2}} \,dv\right)\,\left(\int_0^{+\infty} \dfrac{x^{\frac{n}{2}-1}\,\exp\left\{-\dfrac{x}{2} \right\}}{2^{\frac{n}{2}}\,\Gamma\left(\frac{1}{2}\right)\Gamma\left(\frac{n-1}{2}\right)}\,dx\right)\\
& \left(\int_0^{1} v^{\frac{n-3}{2}}\,\left(1 - v\right)^{-\frac{1}{2}} \,dv\right)\,\left(\int_0^{+\infty} \dfrac{x^{\frac{n}{2}-1}\,\exp\left\{-x \right\}}{\Gamma\left(\frac{1}{2}\right)\Gamma\left(\frac{n-1}{2}\right)}\,dx\right)\\
&= \left(\int_0^{1} v^{\frac{n-3}{2}}\,\left(1 - v\right)^{-\frac{1}{2}} \,dv\right)\,\left( \dfrac{\Gamma\left(\frac{n}{2}\right)}{\Gamma\left(\frac{1}{2}\right)\Gamma\left(\frac{n-1}{2}\right)}\right) = 1.
\end{align*}

On en déduit que
\[ \int_0^{1} v^{\frac{n-3}{2}}\,\left(1 - v\right)^{-\frac{1}{2}} \,dv = \dfrac{\Gamma\left(\frac{1}{2}\right)\Gamma\left(\frac{n-1}{2}\right)}{\Gamma\left(\frac{n}{2}\right)}, \]
et finalement que pour tout \(x\in\mathbb{R}\)
\[f_Y(x) = \left|\begin{array}{ll} \dfrac{x^{\frac{n}{2}-1}}{2^{\frac{n}{2}}\,\Gamma\left(\frac{n}{2}\right)}\,\exp\left\{-\dfrac{x}{2}\right\} & \text{si } x> 0,\\ 0 & \text{sinon;}\end{array}\right.\]
\((\mathcal{P}_n)\) est donc vraie.

\textbf{Conclusion.} La propriété est vraie pour tout
\(n\in\mathbb{N}^\ast\).

\section{Combinaisons linéaires de variables Gaussiennes
indépendantes}\label{II-combinaisons-linuxe9aires-de-variables-gaussiennes-induxe9pendantes-1}

\runinhead{Préliminaires}\phantomsection\addcontentsline{toc}{subsubsection}{Préliminaires}\label{II-answer-CLIGauss-pre}
Cette question a été traitée de manière générale dans le cours. Nous en
proposons une preuve alternative, basée sur le calcul de la fonction de
répartition de la variable aléatoire \(s\,X + m\), qui caractérise sa
loi. Elle dépend clairement des valeurs de \(s\).

\begin{itemize}
\item
  Si \(s = 0\), alors \(s\,X + m\) est toujours égale à \(m\)~: sa loi
  est une masse de Dirac en \(\{m\}\) et pour tout \(x\in\mathbb{R}\) on
  a \(\mathbb{P}\left(s\,X +m \leq x\right) = 1_{[m,+\infty[}(x)\).
\item
  Si \(s\neq 0\), alors pour tout \(x\in\mathbb{R}\) on a
  \[\mathbb{P}\left( s\,X + m \leq x \right) = \left|\begin{array}{ll}\displaystyle\mathbb{P}\left( X \leq \dfrac{x-m}{s} \right) = \int_{-\infty}^{\frac{x-m}{s}} f(u)\,du & \text{si } s>0,\\[1em] \displaystyle \mathbb{P}\left( X \geq \dfrac{x-m}{s} \right) = \int_{\frac{x-m}{s}}^{+\infty} f(u)\,du & \text{si } s<0,\end{array}\right.\]
  qui en posant le changement de variable \(v = s\,x+m\) donne
  \[\mathbb{P}\left( s\,X + m \leq x \right) = \int_{-\infty}^{x} \dfrac{1}{|s|}\,f\left(\dfrac{u-m}{s} \right)\,dx.\]
  Ainsi, \(s\,X+m\) admet une densité, qui pour tout \(x\in\mathbb{R}\)
  est égale à
  \[ \dfrac{1}{|s|}\,f\left(\dfrac{u-m}{s}\right) = \dfrac{1}{\sqrt{2\,\pi}\,|s|}\,\exp\left\{-\dfrac{(x-m)^2}{2\,s^2} \right\}.\]
  On reconnaît la densité d'une loi Normale d'espérance \(m\) et de
  variance \(s^2\).
\end{itemize}

\runinhead{Combinaisons
linéaires}\phantomsection\addcontentsline{toc}{subsubsection}{Combinaisons
linéaires}\label{II-answer-CLIGauss-cl}

\begin{enumerate}
\def\labelenumi{\arabic{enumi}.}
\setcounter{enumi}{1}
\tightlist
\item
  Commençons par supposer que pour tout \(n\in\mathbb{N}^\ast\),
  \(a\in\mathbb{R}^n\) est tel qu'aucune de ses composantes n'est nulle.
  Nous allons montrer par récurrence sur \(n\) que \(S_n^a\) suit une
  loi Normale d'espérance nulle et de variance \(\sum_{i = 1}^n a_i^2\).
  On note cette propriété \((\mathcal{P}_n)\).
\end{enumerate}

\textbf{Initialisation.}

\begin{itemize}
\item
  Si \(n = 1\) et \(a_1 \neq 0\), alors \(S_1^a = a_1\,X_1\) suit une
  loi Normale centrée de variance \(a_1^2\) d'après la question 1;
  \((\mathcal{P}_1)\) est donc vraie.
\item
  Si \(n=2\) et \(a_1,a_2 \neq 0\), alors d'après le cours
  \(S_2^a = a_1\,X_1 + a_2\,X_2\) admet une densité, notée \(f_n^a\),
  égale au produit de convolution des densités \(f_1\) de \(a_1\,X_1\)
  et \(f_2\) de \(a_2\,X_2\). En outre, d'après la question 1, pour tout
  \(x\in\mathbb{R}\) on a
  \(f_1(x) = \dfrac{1}{|a_1|}\,f\left(\dfrac{x}{a_1} \right)\) et
  \(f_2(x) = \dfrac{1}{|a_2|}\,f\left(\dfrac{x}{a_2} \right)\). Par
  conséquent, pour tout \(x\in\mathbb{R}\), \begin{align*}
  f_2^a(x) &= \int_{\mathbb{R}} f_1(x-u)\,f_2(u)\,du = \int_\mathbb{R}\dfrac{1}{|a_1|\,|a_2|}\,f\left(\dfrac{x - u}{a_1}\right)\,f\left(\dfrac{u}{a_2}\right)\,du\\
  &= \dfrac{1}{\sqrt{a_1^2 +a_2^2}}\,f\left(\dfrac{x}{\sqrt{a_1^2+a_2^2}}\right)\, \int_{\mathbb{R}} \dfrac{\sqrt{a_1^2+a_2^2}}{|a_1|\,|a_2|}\,\dfrac{f\left(\dfrac{x - u}{a_1}\right)}{f\left(\dfrac{x}{\sqrt{a_1^2+a_2^2}}\right)}\,f\left(\dfrac{u}{a_2}\right)\,du.
  \end{align*} Or pour tout \(x,u\in\mathbb{R}\) on a \begin{align*}
  \dfrac{f\left(\dfrac{x - u}{a_1}\right)}{f\left(\dfrac{x}{\sqrt{a_1^2+a_2^2}}\right)}\,f\left(\dfrac{u}{a_2}\right) & = \dfrac{1}{\sqrt{2\pi}}\,\dfrac{\exp\left\{-\dfrac{(x-u)^2}{2\,a_1^2} \right\}}{\exp\left\{ -\dfrac{x^2}{2\,(a_1^2+a_2^2)} \right\}}\,\exp\left\{ -\dfrac{u^2}{2\,a_2^2} \right\}\\
  & = \dfrac{1}{\sqrt{2\pi}}\,\exp\left\{-\dfrac{1}{2}\,\left( \dfrac{(x-u)^2}{a_1^2} + \dfrac{u^2}{a_2^2} - \dfrac{x^2}{a_1^2+a_2^2} \right) \right\}\\
  &= \dfrac{1}{\sqrt{2\pi}}\,\exp\left\{ -\dfrac{\left(a_2^2\,x - (a_1^2 + a_2^2)\,u\right)^2}{2\,a_1^2\,a_2^2\,(a_1^2+a_2^2)} \right\}\\
  &= \dfrac{1}{\sqrt{2\pi}}\,\exp\left\{-\dfrac{\left(u - \dfrac{a_2^2\,x}{a_1^2 + a_2^2}\right)^2}{2\,\dfrac{a_1^2\,a_2^2}{a_1^2+a_2^2}}\right\},
  \end{align*} qui multiplié par
  \(\dfrac{\sqrt{a_1^2+a_2^2}}{|a_1|\,|a_2|}\) correspond à la densité
  d'une loi Normale d'espérance \(\dfrac{a_2^2\,x}{a_1^2+a_2^2}\) et de
  variance \(\dfrac{a_1^2\,a_2^2}{a_1^2+a_2^2}\). La précédente
  intégrale vaut donc \(1\) et \((\mathcal{P}_2)\) est vraie.
\end{itemize}

\textbf{Héritage.} Soit maintenant \(n\geq 2\), et supposons
\((\mathcal{P}_{n-1})\) vraie. Alors
\(S_n^a = S_{n-1}^{a_{-n}} + a_n\,X_n\), où
\(a_{-n} := (a_1,\dots,a_{n-1})\). Or \(S_{n-1}^{a_{-n}}\) et
\(a_n\,X_n\) sont des variables gaussiennes centrées, de variances
respectives \(\sum_{i = 1}^{n-1} a_i^2\) d'après \((\mathcal{P}_{n-1})\)
et \(a_n^2\) d'après la question 1. Par \((\mathcal{P}_2)\), \(S_n^a\)
suit donc une loi Normale centrée de variance \(\sum_{i = 1}^n a_i^2\).

\textbf{Conclusion.} On en conclut que \((\mathcal{P}_n)\) est vraie
pour tout \(n\in\mathbb{N}^\ast\).

\begin{enumerate}
\def\labelenumi{\arabic{enumi}.}
\setcounter{enumi}{2}
\tightlist
\item
  Calculons cette covariance~: \begin{align*}
  \text{Cov}\left(S_n^a, S_n^b\right) &= \text{Cov}\left(\sum_{i =1}^n a_i\,X_i, \sum_{j = 1}^n b_j\,X_j\right)\\
  &= \sum_{i = 1}^n a_i\,b_i\,\mathbb{V}\left(X_i\right) + \sum_{1\leq i\neq j \leq n} a_i\,b_j\,\text{Cov}\left(X_i,X_j\right).
  \end{align*} Or par hypothèse \(\mathbb{V}\left(X_i\right) = 1\) pour
  tout \(i\in\{1,\dots,n\}\) et par indépendance on a
  \(\text{Cov}\left(X_i,X_j\right) = 0\) pour tous
  \(i,j \in \{1,\dots,n\}\) tels que \(i\neq j\). Ainsi,
  \[ \text{Cov}\left(S_n^a, S_n^b\right) = \sum_{i = 1}^n a_i\,b_i \]
  qui est nulle ssi \(a\) et \(b\) sont orthogonaux.
\end{enumerate}

\section{Loi de Maxwell}\label{II-loi-de-maxwell-1}

\runinhead{Question
1}\phantomsection\addcontentsline{toc}{subsubsection}{Question
1}\label{II-answer-Maxwell-1}
La densité du triplet \((X, Y, Z)\) étant invariante par rotation, elle
est de la forme \(f_{(X,Y,Z)} (x, y, z) = \phi(x^2 + y^2 + z^2)\). De
plus, les variables \(X, Y\) et \(Z\) étant indépendantes, on a
\[f_{(X,Y,Z)} (x, y, z) = f_X (x)f_Y (y)f_Z (z) \; \mbox{ pour tout $(x, y, z) \in \mathbb{R}^3$.} \]
Soit \((x_0,y_0,z_0) \in \mathbb{R}^3\)\,; comme
\(\phi(x_0^2 + y_0^2 + z_0^2) > 0\) on a
\(f_{(X, Y, Z)}(x_0, y_0, z_0) > 0\) et donc \(f_X(x_0) \neq 0\),
\(f_Y(y_0) \neq 0\) et \(f_Z(z_0) \neq 0\). Pour tout
\(x \in \mathbb{R}^3\), comme \[
f_X(x) f_Y(y_0) f_Z(z_0) = \phi(x^2 + y_0^2 + z^2) = \phi(y_0^2 + x^2 + z_0^2)
= f_X(y_0) f_Y(x) f_Z(z_0),
\] on a l'égalité \[
f_X(x) = \frac{f_X(y_0)}{f_Y(y_0)} f_Y(x).
\] Comme on a par ailleurs \[
\int_{-\infty}^{+\infty} f_X(x) \, dx = \int_{-\infty}^{+\infty} f_Y(x) \, dx = 1,
\] il est nécessaire que le ratio \({f_X(y_0)}/{f_Y(y_0)}\) soit égal à
\(1\). On a donc \(f_X = f_Y\)\,; la preuve que \(f_Y = f_Z\) s'obtient
de manière similaire.

Pour finir, puisque pour tout \(x \in \mathbb{R}\) on a
\(f(x)^3 = f_{(X, Y, Z)}(x, x, x) = \phi(3 x^2) > 0\), la fonction \(f\)
est bien positive en tout point.

\runinhead{Question
2}\phantomsection\addcontentsline{toc}{subsubsection}{Question
2}\label{II-answer-Maxwell-2}
L'égalité
\[f(x)f(y)f(z) = \phi(x^2 + y^2 + z^2 )  \; \text{ pour tout } (x, y, z) \in \mathbb{R}^3\]
implique \(f (x)f'(y)f (z) = 2y\phi' (x^2 + y^2 + z^2 )\) et
\(f' (x)f (y)f (z) = 2x\phi' (x^2 + y^2 + z^2 ).\) On en déduit que
\[xf (x)f' (y)f (z) = f' (x)yf (y)f (z)  \; \text{ pour tout } (x, y, z) \in \mathbb{R}^3.\]

Soit \((x, y_0, z_0) \in \mathbb{R}^3\) avec \(y_0 \neq 0\)\,; on a
\(f(x) > 0\), \(f(y_0) > 0\) et \(f(z_0) > 0\). En posant \[
c = \frac{f'(y_0) f(z_0)}{y_0 f(y_0) f(z_0)},
\] on voit que
\[f' (x) = cxf(x) \; \mbox{pour tout $x \in \mathbb{R}$.}\] On peut
réécrire cette équation différentielle sous la forme \[
(\ln f(x))' = cx.
\] En intégrant les deux membres de l'équation entre \(0\) et \(x\), on
obtient \[
\ln f(x) - \ln f(0) = c\frac{x^2}{2}
\] et donc \[
f(x) = f(0) e^{c x^2/2},
\] ce qui est bien la forme cherchée \(f(x) = a e^{cx^2/2}\).

\runinhead{Question
3}\phantomsection\addcontentsline{toc}{subsubsection}{Question
3}\label{II-answer-Maxwell-3}
Comme \(\int_\mathbb{R}f (x) dx = 1\), on a nécessairement \(c < 0\) et
on en déduit que
\[f (x) = \frac{1}{\sqrt{2\pi \sigma^2}}\exp\left(-\frac{x^2}{2\sigma^2}\right).\]
Les variables aléatoires \(X\), \(Y\) et \(Z\) suivent donc loi
gaussienne centrée et la covariance du vecteur est \(\sigma^2 I_3\), où
\(I_3\) est la matrice identité.

\runinhead{Question
4}\phantomsection\addcontentsline{toc}{subsubsection}{Question
4}\label{II-answer-Maxwell-4}
On a
\[\frac{1}{2} m\mathbb{E}(\|V \|^2 ) = \frac{1}{2} m\mathbb{E}(V_1^2 +V_2^2 + V_3^2).\]
Comme \(\mathbb{E}(V_i^2) = \sigma^2\), on en déduit que
\(E_c = \frac{3}{2}m \sigma^2\) et par conséquent que
\[\sigma^2 = \frac{kT}{m}.\]

\runinhead{Question
5}\phantomsection\addcontentsline{toc}{subsubsection}{Question
5}\label{II-answer-Maxwell-5}
À l'aide de la méthode de la fonction muette, on montre que \(V_1^2\)
suit une loi gamma de paramètres \((1/2,\frac{1}{2\sigma^2})\). En
effet, en appliquant le changement de variable \(y=x^2\) sur
\(\mathbb{R}_-^\ast\) et \(\mathbb{R}_+^\ast\), on obtient~:
\[f_{V_1^2} (y) = \left(f(-\sqrt{y}) + f(\sqrt{y})\right) \frac{1}{2\sqrt{y}} 1_{]0,+\infty[}(x) = f(\sqrt{y}) \frac{1}{\sqrt{y}} 1_{]0,+\infty[}(x)\]
par parité de \(f\), soit \begin{align*}
f_{V_1^2} (y) &= \frac{1}{\sqrt{\pi}}\left(\frac{1}{2\sigma^2}\right)^{1/2} y^{-1/2} \exp\left(-\frac{y}{2\sigma^2}\right)\\
              &= \frac{1}{\Gamma(1/2)}\left(\frac{1}{2\sigma^2}\right)^{1/2} y^{-1/2} \exp\left(-\frac{y}{2\sigma^2}\right)
\end{align*} On déduit de l'indication que
\(\|V\|^2 = V_1^2 +V_2^2 + V_3^2\) suit une loi gamma de paramètres
\((3/2,\frac{1}{2\sigma^2})\). Sa densité est
\[f_{\|V\|^2}(x) = \frac{1}{\sqrt{2\pi}\sigma^3}\sqrt{x} e^{- x/2\sigma^2}1_{]0,+\infty[}(x).\]

À l'aide de la méthode de la fonction muette, en appliquant le
changement de variable \(v = \sqrt{x}\) sur \(\mathbb{R}_+^\ast\) et
avec \(\sigma^2 = \frac{kT}{m}\), on montre enfin que \(\|V\|\) est une
variable de densité~:

\[f_{\|V\|}(v) = \frac{\sqrt{2}}{\sqrt{\pi}}\left(\frac{m}{kT}\right)^{3/2}v^2 e^{- mv^2/2kT} 1_{]0,+\infty[}(v).\]

Il s'agit de la densité de la loi de Maxwell.
